
\section{Native Sparse Attention: Sliding-Window Local Context with Dynamic Top-K Global Routing}


\textbf{Authors:}
\begin{itemize}
\item {{AUTHOR\_1\_NAME}}  -- {{AUTHOR\_1\_EMAIL}}
\end{itemize}
\begin{itemize}
\item {{AUTHOR\_2\_NAME}}  -- {{AUTHOR\_2\_EMAIL}}
\end{itemize}

\textbf{Affiliations:}
\begin{itemize}
\item   {{UNIVERSITY\_1}}, {{DEPARTMENT\_1}}
\end{itemize}
\begin{itemize}
\item   {{UNIVERSITY\_2}}, {{DEPARTMENT\_2}}
\end{itemize}

\textbf{Corresponding author:} {{CORRESPONDING\_AUTHOR}}
\textbf{Preprint date:} {{DATE}}
\textbf{Version:} v{{VERSION}}
\textbf{Code:} https://github.com/{{ORG}}/SVD\_Project/tree/main/framework/attentions/nsa.py
\textbf{Model weights:} https://huggingface.co/{{HF\_ORG}}
\textbf{License:} {{LICENSE}}



\subsection{Abstract}

Sparse attention methods exploit the observation that most attention mass lies 
within a small neighborhood around each query. Fixed-window methods 
(Longformer, BigBird) capture local context cheaply but miss distant 
dependencies. Content-based methods (Routing Transformer, Reformer) capture 
distant dependencies but add hashing or clustering overhead. We present 
\textbf{Native Sparse Attention} (NSA), which combines \textbf{both}: a fixed sliding 
window captures local context, while a \textbf{dynamic top-$k$ routing} selects 
the $k$ most relevant distant keys per query based on actual attention scores. 
The selection is \textbf{fully vectorized} -- no Python loop, no clustering, no 
hashing -- making NSA easy to implement, train, and reason about.

NSA is parameterized by two scalars: window size $w$ and top-$k$ count $k$. 
For each query $i$, the retained keys are:

\[
\mathcal{S}_i = \underbrace{\{j : |i - j| \leq w\}}_{\text{local window}} \cup \underbrace{\text{top-}k_{j' \notin \mathcal{W}_i} s_{ij'}}_{\text{dynamic global}}
\]

The union is computed via a boolean mask that is applied before softmax. 
Padding, causality, and GQA are all preserved.

We evaluate NSA on DistilBERT, RoBERTa-base, and GPT-2 on SST-2 with 
$w = {{WINDOW_SIZE}}$, $k = {{TOP_K}}$. NSA achieves \textbf{{{D\_ACC}}\%}, 
\textbf{{{R\_ACC}}\%}, and \textbf{{{G\_ACC}}\%} accuracy, within \textbf{{{D\_DELTA}}\%}, 
\textbf{{{R\_DELTA}}\%}, and \textbf{{{G\_DELTA}}\%} of the dense baselines. Attention 
density drops to \textbf{{{DENSITY}}\%} (from 100\% dense). The dynamic top-$k$ 
selection shifts systematically with layer depth, revealing head-level 
content specialization.

\textbf{Keywords:} sparse attention, sliding window, dynamic top-k, content-aware 
routing, efficient transformer



\subsection{Lay Summary}

Most language-model attention weight lies on nearby tokens. Some attention 
heads, however, need to look far away. Existing methods handle one regime or 
the other: fixed windows see only nearby tokens; content-based methods need 
expensive clustering or hashing.

Our method NSA runs both at once: for every token, we keep a fixed window 
of nearby context AND dynamically pick the few most-relevant distant tokens 
based on their actual attention scores. The selection is done in a single 
vectorized operation -- no loop, no clustering.

On three pretrained models, NSA retains most accuracy while cutting attention 
density to about {{DENSITY}}\%. The set of dynamically-selected tokens changes 
systematically with layer depth, revealing what each layer cares about.



\subsection{Table of Contents}

1. [Introduction](01\_intro.md)
2. [Related Work](02\_related.md)
3. [Background](03\_background.md)
4. [Method](04\_method.md)
5. [Architecture Integration](05\_architecture.md)
6. [Experimental Setup](06\_setup.md)
7. [Results](07\_results.md)
8. [Analysis](08\_analysis.md)
9. [Discussion](09\_discussion.md)
10. [Limitations](10\_limitations.md)
11. [Conclusion](11\_conclusion.md)
12. [Reproducibility](12\_reproducibility.md)
13. [References](13\_references.md)

\textbf{Appendices:}
\begin{itemize}
\item [A. Mathematical Derivations](A\_math.md)
\end{itemize}
\begin{itemize}
\item [B. Implementation Details](B\_implementation.md)
\end{itemize}
\begin{itemize}
\item [C. Full Experimental Data](C\_data.md)
\end{itemize}
\begin{itemize}
\item [D. Negative Results](D\_negative\_results.md)
\end{itemize}
\begin{itemize}
\item [E. Ethical Considerations](E\_ethics.md)
\end{itemize}